%==============================================================
\section{在物理学中的应用}
%==============================================================

\subsection{量子力学中的对称性}

\subsubsection{角动量与 $\mathfrak{so}(3)$}

\textbf{角动量算子满足 $\mathfrak{so}(3)$ 代数：}
\[
[\hat{L}_i, \hat{L}_j] = i\hbar\varepsilon_{ijk}\hat{L}_k
\]

\textbf{不可约表示 → 角动量量子化：}
\begin{itemize}
  \item $\hat{L}^2|l,m\rangle = \hbar^2 l(l+1)|l,m\rangle$
  \item $\hat{L}_z|l,m\rangle = \hbar m|l,m\rangle$
  \item $l = 0, 1, 2, \ldots$；$m = -l, \ldots, l$
\end{itemize}

\textbf{球谐函数 $Y_l^m$ 是 $SO(3)$ 的第 $l$ 个不可约表示的基函数。}

\subsubsection{自旋与 $\mathfrak{su}(2)$}

\textbf{自旋算子满足 $\mathfrak{su}(2)$ 代数：}
\[
[\hat{S}_i, \hat{S}_j] = i\hbar\varepsilon_{ijk}\hat{S}_k
\]

电子自旋 $s = 1/2$：$SU(2)$ 的 2 维表示。

\textbf{为什么是 $SU(2)$ 而非 $SO(3)$？} 因为 $\pi_1(SO(3)) = \mathbb{Z}_2$，$SO(3)$ 的投影表示（允许相位）等价于 $SU(2)$ 的普通表示。

\subsubsection{氢原子与 $SO(4)$}

氢原子有隐藏对称性：除了角动量 $\mathbf{L}$，还有 Runge--Lenz 向量 $\mathbf{A}$。

$\mathbf{L}$ 和 $\mathbf{A}$ 共同生成 $\mathfrak{so}(4) \cong \mathfrak{su}(2) \oplus \mathfrak{su}(2)$。

\textbf{结果：} 能级 $E_n$ 的简并度是 $n^2$（而非仅 $2l+1$），这由 $SO(4)$ 的不可约表示解释。

\subsection{规范场论}

\subsubsection{规范对称性 → 力}

\textbf{核心思想：} 要求拉格朗日量在\textbf{局部}规范变换下不变 → 必须引入规范场（联络）→ 规范场就是力。

\begin{center}
\begin{tabular}{p{2cm}p{3cm}p{3cm}p{4cm}}
\toprule
\textbf{规范群} & \textbf{规范场} & \textbf{力} & \textbf{粒子} \\
\midrule
$U(1)$ & $A_\mu$（1 个） & 电磁力 & 光子 \\
$SU(2)$ & $W_\mu^a$（3 个） & 弱力 & $W^\pm, Z^0$ \\
$SU(3)$ & $G_\mu^a$（8 个） & 强力 & 8 个胶子 \\
\bottomrule
\end{tabular}
\end{center}

\subsubsection{标准模型的群结构}

\[
G_{SM} = SU(3)_C \times SU(2)_L \times U(1)_Y
\]

\begin{itemize}
  \item $SU(3)_C$：色对称（强相互作用）
  \item $SU(2)_L$：弱同位旋（弱相互作用）
  \item $U(1)_Y$：弱超荷
\end{itemize}

\textbf{电弱对称性破缺：} $SU(2)_L \times U(1)_Y \xrightarrow{\text{Higgs}} U(1)_{EM}$

\subsubsection{规范场的数学结构}

\textbf{规范场 = 主丛上的联络：}
\[
A = A_\mu^a T_a\, dx^\mu
\]

\textbf{场强 = 曲率：}
\[
F = dA + A \wedge A = \frac{1}{2}F_{\mu\nu}^a T_a\, dx^\mu \wedge dx^\nu
\]

\textbf{Yang--Mills 作用量：}
\[
S_{YM} = -\frac{1}{2}\int \text{tr}(F \wedge *F) = -\frac{1}{4}\int F_{\mu\nu}^a F^{a\mu\nu}\,d^4x
\]

\subsection{广义相对论与时空对称性}

\textbf{时空对称群：}
\begin{itemize}
  \item 牛顿力学：Galilei 群（10 维）
  \item 狭义相对论：Poincaré 群 $ISO(1,3)$（10 维）
  \item 广义相对论：微分同胚群 $\text{Diff}(M)$（无穷维）
\end{itemize}

\textbf{Poincaré 群 = Lorentz 群 + 平移：}
\[
ISO(1,3) = SO(1,3) \ltimes \mathbb{R}^4
\]

\textbf{Wigner 分类：} 粒子 = Poincaré 群的不可约酉表示，由 $(m, s)$ 标记（质量、自旋）。

\subsection{凝聚态物理}

\begin{center}
\begin{tabular}{p{4cm}p{4cm}p{4cm}}
\toprule
\textbf{对称性} & \textbf{破缺} & \textbf{结果} \\
\midrule
$U(1)$（相位） & 超导 & 磁通量子化、Josephson 效应 \\
$SO(3)$（自旋） & 铁磁 & 磁畴、自旋波 \\
平移对称性 & 晶体 & 声子、能带结构 \\
$SU(2) \times SU(2)$ & 手征对称性破缺 & 核子质量 \\
\bottomrule
\end{tabular}
\end{center}

\textbf{Goldstone 定理：} 连续对称性自发破缺 → 无质量 Goldstone 玻色子。

\textbf{Higgs 机制：} 规范对称性自发破缺 → 规范玻色子获得质量（Goldstone 玻色子被"吃掉"）。

%==============================================================
\section{在微分方程中的应用}
%==============================================================

\subsection{对称性方法求解 ODE}

（详见微分方程专题第四层，此处总结要点）

\begin{center}
\begin{tabular}{p{4cm}p{4cm}p{4cm}}
\toprule
\textbf{对称性} & \textbf{生成元} & \textbf{效果} \\
\midrule
$x$ 平移 & $\partial_x$ & 降阶（缺 $x$） \\
$y$ 平移 & $\partial_y$ & 降阶（缺 $y$） \\
尺度变换 & $x\partial_x + y\partial_y$ & 齐次方程降阶 \\
旋转 & $-y\partial_x + x\partial_y$ & 极坐标简化 \\
\bottomrule
\end{tabular}
\end{center}

\subsection{对称性方法求解 PDE}

\textbf{热方程 $u_t = ku_{xx}$ 的对称群（6 维）：}
\begin{itemize}
  \item $x$ 平移：$\partial_x$
  \item $t$ 平移：$\partial_t$
  \item $u$ 的线性叠加：$u\partial_u$
  \item 尺度：$(2t\partial_t + x\partial_x)$
  \item Galilei boost：$(2kt\partial_x + xu\partial_u)$
  \item 特殊共形：$(4k^2t^2\partial_t + 4ktx\partial_x + (2kt - x^2)u\partial_u)$
\end{itemize}

\textbf{利用对称性构造特殊解：}
\begin{itemize}
  \item 平移不变 → 行波解 $u = f(x - ct)$
  \item 尺度不变 → 自相似解 $u = t^{-\alpha}f(x/\sqrt{t})$（热核！）
  \item 旋转不变 → 径向解 $u = u(r, t)$
\end{itemize}

\subsection{守恒律与矩映射}

\textbf{Noether 定理的李群表述：}

若作用量 $S = \int L\,dt$ 在一参数群 $g_\varepsilon$ 下不变，则：
\[
Q = \frac{\partial L}{\partial \dot{q}} \cdot \delta q
\]
是守恒量。

\textbf{矩映射（Moment Map）：} 将对称群的李代数映射到相空间上的守恒量：
\[
\mu: \mathfrak{g} \to C^\infty(M), \quad X \mapsto \mu_X
\]

满足 $\{\mu_X, f\} = X_M(f)$（$X_M$ 是 $X$ 在相空间 $M$ 上诱导的向量场）。

\textbf{例子：}
\begin{itemize}
  \item 平移对称 → 动量 $p$
  \item 旋转对称 → 角动量 $\mathbf{L} = \mathbf{r} \times \mathbf{p}$
  \item 时间平移 → 能量 $H$
\end{itemize}

%==============================================================
\section{在计算机图形学中的应用}
%==============================================================

\subsection{三维旋转与四元数}

\subsubsection{旋转的多种表示}

\begin{center}
\begin{tabular}{p{3cm}p{3cm}p{3cm}p{3cm}}
\toprule
\textbf{表示} & \textbf{参数数} & \textbf{奇异性} & \textbf{应用} \\
\midrule
旋转矩阵 $R \in SO(3)$ & 9（6 约束） & 无 & 通用变换 \\
欧拉角 & 3 & 有（万向锁） & 动画关键帧 \\
轴角 $(\hat{n}, \theta)$ & 3 & $\theta = 0$ 处 & 物理模拟 \\
四元数 $q \in SU(2)$ & 4（1 约束） & 无 & 插值、组合 \\
\bottomrule
\end{tabular}
\end{center}

\subsubsection{四元数与 $SU(2)$}

单位四元数 $q = \cos\frac{\theta}{2} + \sin\frac{\theta}{2}(\hat{n}_x i + \hat{n}_y j + \hat{n}_z k)$

\textbf{旋转作用：} $\mathbf{v}' = q\mathbf{v}q^{-1}$（将 $\mathbf{v}$ 视为纯虚四元数）

\textbf{优势：}
\begin{itemize}
  \item 无万向锁
  \item 组合旋转 = 四元数乘法（$O(1)$ vs 矩阵乘法 $O(27)$）
  \item 插值自然（Slerp）
  \item 归一化便宜（4 个分量）
\end{itemize}

\subsubsection{Slerp（球面线性插值）}

\[
\text{Slerp}(q_0, q_1, t) = \frac{\sin((1-t)\Omega)}{\sin\Omega}q_0 + \frac{\sin(t\Omega)}{\sin\Omega}q_1
\]

其中 $\cos\Omega = q_0 \cdot q_1$（四元数点积）。

\textbf{几何意义：} 在 $S^3$（单位四元数空间）上沿大圆弧插值。

\subsection{刚体运动与 $SE(3)$}

\subsubsection{刚体变换}

\[
T = \begin{pmatrix} R & \mathbf{t} \\ \mathbf{0}^T & 1 \end{pmatrix} \in SE(3)
\]

\subsubsection{$\mathfrak{se}(3)$ 与螺旋运动}

\[
\mathfrak{se}(3) = \left\{\begin{pmatrix} [\omega]_\times & \mathbf{v} \\ \mathbf{0}^T & 0 \end{pmatrix} : \omega, \mathbf{v} \in \mathbb{R}^3\right\}
\]

\textbf{指数映射（螺旋运动）：}
\[
\exp\begin{pmatrix} [\omega]_\times & \mathbf{v} \\ \mathbf{0}^T & 0 \end{pmatrix} = \begin{pmatrix} e^{[\omega]_\times} & V\mathbf{v} \\ \mathbf{0}^T & 1 \end{pmatrix}
\]

其中 $V = I + \frac{1-\cos\theta}{\theta^2}[\omega]_\times + \frac{\theta - \sin\theta}{\theta^3}[\omega]_\times^2$。

\textbf{在 CG 中：} 骨骼动画的关节变换、相机运动、机器人学。

\subsection{球谐函数与 $SO(3)$ 表示}

\subsubsection{球谐函数 = $SO(3)$ 的不可约表示基}

$Y_l^m(\theta, \phi)$ 是 $SO(3)$ 的第 $l$ 个不可约表示（维度 $2l+1$）的基函数。

\textbf{旋转下的变换：}
\[
Y_l^m(R^{-1}\hat{n}) = \sum_{m'=-l}^{l} D^l_{m'm}(R)\, Y_l^{m'}(\hat{n})
\]

$D^l_{m'm}(R)$ 是 Wigner $D$-矩阵（$(2l+1) \times (2l+1)$ 酉矩阵）。

\subsubsection{在渲染中的应用}

\textbf{环境光照：} 用球谐系数表示：
\[
L(\hat{n}) \approx \sum_{l=0}^{L}\sum_{m=-l}^{l} c_l^m Y_l^m(\hat{n})
\]

通常 $L = 2$（9 个系数）即可表示低频环境光。

\textbf{PRT（预计算辐射传输）：} 将传输算子在球谐基上对角化，实现实时全局光照。

\textbf{球面卷积：} 在球谐域变成逐元素乘法：
\[
\widehat{f * g}_l^m = \hat{f}_l \cdot \hat{g}_l^m
\]

（卷积定理的球面版本）

\subsection{等变神经网络}

\textbf{核心思想：} 网络层在群作用下等变：
\[
f(g \cdot x) = g \cdot f(x)
\]

\textbf{实现：} 用不可约表示的张量积构造等变层。

\textbf{应用：}
\begin{itemize}
  \item 分子性质预测（$SE(3)$ 等变）
  \item 流体模拟（旋转等变）
  \item 3D 点云处理
\end{itemize}

\subsection{微分几何与网格处理}

\subsubsection{曲面上的对称性}

\textbf{等距变换群：} 保持第一基本形式不变的变换。

\textbf{共形变换群：} 保持角度不变的变换（无穷维！）。

\textbf{在 CG 中：}
\begin{itemize}
  \item 纹理映射：共形映射（保持角度）
  \item 网格参数化：等距映射（保持距离）
  \item 形状匹配：在对称群下不变的特征
\end{itemize}

\subsubsection{离散微分几何中的群结构}

\textbf{离散联络：} 网格边上的旋转（$SO(3)$ 元素），描述向量沿边的平行移动。

\textbf{离散曲率：} 绕顶点一圈的旋转角度亏损 = 高斯曲率的离散版本。

\subsection{拓扑在 CG 中的应用}

\subsubsection{网格拓扑}

\textbf{Euler 公式：} $V - E + F = 2 - 2g$（$g$ 是亏格）

\begin{center}
\begin{tabular}{p{3cm}p{2cm}p{5cm}}
\toprule
\textbf{拓扑} & \textbf{亏格} & \textbf{例子} \\
\midrule
球面 & 0 & 封闭物体 \\
环面 & 1 & 甜甜圈、手柄 \\
双环面 & 2 & 眼镜框 \\
\bottomrule
\end{tabular}
\end{center}

\textbf{影响：}
\begin{itemize}
  \item 参数化：球面拓扑可以展开到平面（需要切割），环面可以无切割展开
  \item 向量场：球面上不存在处处非零的连续切向量场（Hairy Ball 定理）
  \item 网格生成：拓扑决定了网格的连接方式
\end{itemize}

\subsubsection{Hairy Ball 定理}

$S^2$ 上不存在处处非零的连续切向量场。

\textbf{在 CG 中：}
\begin{itemize}
  \item 球面上的纹理坐标必有奇点（至少一个极点）
  \item 风向/毛发方向场在球面上必有零点
  \item 这解释了为什么地球仪的地图投影必有畸变
\end{itemize}

\subsubsection{同调与网格分析}

\textbf{Betti 数：}
\begin{itemize}
  \item $b_0$：连通分量数
  \item $b_1$：独立环路数（"洞"的个数）
  \item $b_2$：封闭腔体数
\end{itemize}

\textbf{在 CG 中：} 网格质量检查、拓扑简化、特征提取。

%==============================================================
\section{总结：李群--李代数--拓扑的统一视角}
%==============================================================

\subsection{三条主线的交汇}

\[
\boxed{
\begin{aligned}
&\text{李群（全局对称性）} \xleftrightarrow{\text{指数映射}} \text{李代数（无穷小对称性）} \\
&\text{李群} \xleftrightarrow{\text{基本群/覆盖}} \text{拓扑（全局结构）} \\
&\text{李代数} \xleftrightarrow{\text{表示论}} \text{物理（粒子/量子数）}
\end{aligned}
}
\]

\subsection{核心对应表}

\begin{center}
\begin{tabular}{p{3cm}p{3cm}p{3cm}p{3cm}}
\toprule
\textbf{数学} & \textbf{物理} & \textbf{CG} & \textbf{微分方程} \\
\midrule
李群 & 对称群 & 变换群 & 对称性 \\
李代数 & 守恒量代数 & 无穷小变换 & 降阶生成元 \\
表示 & 粒子分类 & 球谐/傅里叶 & 特征函数 \\
指数映射 & 时间演化 & 旋转插值 & 矩阵指数解 \\
覆盖群 & 自旋 & 四元数 & 多值解 \\
特征类 & 拓扑荷 & 网格拓扑 & 指标定理 \\
\bottomrule
\end{tabular}
\end{center}

\subsection{学习路径建议}

\begin{enumerate}
  \item \textbf{入门：} $SO(2)$、$SO(3)$、四元数、旋转矩阵（CG 直接可用）
  \item \textbf{进阶：} 李代数、指数映射、BCH 公式（物理模拟、机器人学）
  \item \textbf{深入：} 表示论、球谐函数（渲染、量子力学）
  \item \textbf{抽象：} 纤维丛、特征类、同伦群（规范场论、拓扑物态）
\end{enumerate}

\subsection{推荐参考}

\begin{itemize}
  \item Hall, \textit{Lie Groups, Lie Algebras, and Representations}（数学系标准教材）
  \item Stillwell, \textit{Naive Lie Theory}（直观入门）
  \item Georgi, \textit{Lie Algebras in Particle Physics}（物理应用）
  \item Nakahara, \textit{Geometry, Topology and Physics}（拓扑与物理）
  \item Crane, \textit{Discrete Differential Geometry}（CG 中的离散几何）
  \item Sola et al., \textit{A micro Lie theory for state estimation in robotics}（工程应用）
\end{itemize}